\documentclass[11pt]{article}
\usepackage[margin=1in]{geometry}
\usepackage{amsmath,amssymb,amsthm}
\usepackage{booktabs}
\usepackage{hyperref}
\usepackage{xcolor}
\usepackage{listings}
\usepackage{graphicx}

\title{Independent Replication Report:\\
Spectral Discretization of Darcy's Equations Coupled with the Heat Equation\\
{\large Bernardi, Maarouf, Yakoubi (IMA J.\ Numer.\ Anal., 2016)}}
\author{PDE-100 Replication Project}
\date{2026-07-01}

\begin{document}
\maketitle

\begin{abstract}
\noindent\textbf{Verdict: REPLICATED.} An independent, from-scratch Legendre--Gauss--Lobatto (GLL)
spectral Galerkin solver for the coupled Darcy + heat system reproduces every reproducible central
numerical claim of the accuracy test in Sec.~5.2 (Figs.~1--2) of Bernardi, Maarouf, Yakoubi
(2016): spectral (exponential) convergence in $N$ of the $L^2(u,p,T)$ and $H^1(p,T)$ errors,
the machine-precision floor near $N\!\approx\!20$, indistinguishability of exact and discrete
solutions at $N=17$, and convergence of the decoupled fixed-point iteration (Thm.~5.1). Two
independent free LLM judges concur (\textsc{replicated}). A latent typo in the printed source
terms of eq.~(5.6) was diagnosed symbolically; it does not affect the paper's substance.
\end{abstract}

\section{Paper summary}
\label{sec:paper}
The authors analyze the temperature field of a fluid moving through a porous medium under a
Boussinesq / quasi-stationary approximation, with a nonlinear buoyancy source $F(T)$ modeling an
exothermic reaction. On $\Omega=(-1,1)^d$,
\begin{align}
  \alpha\,\mathbf{u} + \nabla p &= F(T) &&\text{in }\Omega,\\
  \nabla\!\cdot\!\mathbf{u} &= 0        &&\text{in }\Omega,\\
  -\lambda\,\Delta T + (\mathbf{u}\!\cdot\!\nabla)T &= h &&\text{in }\Omega,
\end{align}
with $\mathbf{u}\!\cdot\!\mathbf{n}=0$ on $\partial\Omega$, $T=T_\star$ on $\Gamma_\star$,
$\partial T/\partial n=\theta_\sharp$ on $\Gamma_\sharp$.

\paragraph{Contributions.}
(i) Existence/uniqueness of a weak solution; (ii) a GLL spectral $\mathbb{P}_N\times\mathbb{P}_N
\times\mathbb{P}_N$ Galerkin discretization with numerical integration (Sec.~3); (iii) an optimal
a priori estimate (Thm.~4.7, eq.~4.16) of the form
\[
  \|u-u_N\| + \|p-p_N\|_{H^1} + \|\theta-\theta_N\|_{H^1}
  \;\le\; c\,N^{d/6-s}\|u\|_{H^s} + c\,N^{-s}\bigl(\|p\|_{H^{s+1}}+\|\theta\|_{H^{s+1}}\bigr)+\cdots,
\]
i.e.\ error decays faster than any algebraic power of $N$ for analytic solutions --- spectral
convergence; (iv) a decoupled fixed-point iteration (5.1--5.3) with a proven contraction
(Thm.~5.1); (v) numerical experiments in FreeFEM3D (Sec.~5).

\paragraph{Accuracy test (Sec.~5.2, eq.~5.6).}
Manufactured analytic solution on $\Omega=(-1,1)^2$:
$u_1=-\sin(\pi x)\cos(\pi y)$, $u_2=\cos(\pi x)\sin(\pi y)$, $p=-\sin(\pi x)\cos(\pi y)$,
$T=(1/\pi^2)\cos(\pi x)\sin(\pi y)$, permeability $\alpha(x,y)=1/(T^2+1)$ (interpolated $I_N\alpha$),
$\lambda=1$. The paper's Fig.~1 shows spectral decay of $L^2$ and $H^1$ errors as $N$ runs
$5\to 25$; beyond $N\!\approx\!20$ machine precision dominates. Fig.~2 shows exact vs.\ discrete $T$
at $N=17$: visually indistinguishable.

\section{Claims table}
\begin{center}
\small
\begin{tabular}{clll}
\toprule
ID & Claim & Testable & Outcome \\
\midrule
C1 & $L^2$ errors of $u,p,T$ decay spectrally & yes & Reproduced ($e^{-bN}$, $b\!\approx\!1.9\text{--}2.2$) \\
C2 & $H^1$ errors of $p,T$ decay spectrally  & yes & Reproduced ($b\!\approx\!2.0\text{--}2.05$) \\
C3 & Machine-precision floor beyond $N\!\approx\!20$ & yes & Reproduced (floor at $N\!\approx\!16\text{--}18$) \\
C4 & Exact vs discrete indistinguishable at $N=17$ (Fig.~2) & yes & Reproduced ($\max|T-T_N|=1.5\text{e-}14$) \\
C5 & Decoupled fixed-point (5.1--5.3) converges (Thm 5.1) & yes & Reproduced (iters $4\to 1$) \\
C6 & Optimal a priori estimate (4.16) & partial & Consistent (proof not re-derived) \\
C7 & Existence/uniqueness proofs & no & Out of scope \\
C8 & Horton--Rogers--Lapwood physical case (Figs 3--5) & yes & Not attempted (no error metric) \\
\bottomrule
\end{tabular}
\end{center}

\section{Independent method}
\label{sec:method}
\paragraph{Data source.} OA preprint from HAL \texttt{hal-01085011}
(MD5 \texttt{2d6ead2ce797287b0718d8cc156a1ecd}, 23~pp.). No public author code exists (paper used
an internal FreeFEM3D spectral code) --- this is a from-scratch reimplementation, not a rerun.

\paragraph{Tools.} Python 3.14.6, numpy 2.4.3, scipy 1.18.0, sympy 1.14.0, matplotlib. macOS
(CherryRd). Light compute --- local, no GPU.

\paragraph{Discretization (faithful to Sec.~3).}
Tensor-product GLL grid $(N{+}1)^2$ on $(-1,1)^2$; discrete inner product $(\cdot,\cdot)_N$ =
diagonal GLL mass. \emph{Darcy block:} eliminate $\mathbf{u}$ pointwise from
$\alpha\mathbf{u}=F-\nabla p$; impose $(\mathbf{u},\nabla q)_N=0$ to obtain a symmetric weak-Darcy
system for $p$ (mean-zero constraint via Lagrange multiplier); recover $\mathbf{u}$ pointwise.
\emph{Heat block:} $(\nabla\theta,\nabla\varphi)_N + ((\mathbf{u}\!\cdot\!\nabla)\theta,\varphi)_N
= (h,\varphi)_N$ with Dirichlet data from the MMS. \emph{Coupling:} paper's decoupled fixed-point
scheme (5.1--5.3).

\paragraph{Steps.}
\texttt{spectral\_gll.py} (self-tested to $\sim$1e-14) $\to$ \texttt{verify\_mms.py} $\to$
\texttt{darcy\_heat\_solver.py} (sweep $N=5..25$) $\to$ \texttt{analyze.py} (exponential fits) $\to$
\texttt{fig\_solution.py} (Fig.~2 replica) $\to$ \texttt{judge.py} (LLM judges).

\section{Results vs.\ paper}
\label{sec:results}
\begin{center}
\small
\begin{tabular}{rlllll r}
\toprule
$N$ & $L^2(u)$ & $L^2(p)$ & $L^2(T)$ & $H^1(p)$ & $H^1(T)$ & fp iters \\
\midrule
 5 & 4.44e-01 & 2.00e-02 & 1.80e-03 & 1.31e-01 & 1.12e-02 & 4 \\
 8 & 5.70e-03 & 1.60e-04 & 1.04e-05 & 1.47e-03 & 7.96e-05 & 4 \\
11 & 2.35e-05 & 3.23e-07 & 1.68e-08 & 4.08e-06 & 2.02e-07 & 3 \\
14 & 4.34e-08 & 4.31e-10 & 2.21e-11 & 6.85e-09 & 3.22e-10 & 2 \\
17 & 4.21e-11 & 2.58e-13 & 1.24e-14 & 5.00e-12 & 2.26e-13 & 1 \\
20 & 3.50e-14 & 4.94e-15 & 2.84e-16 & 2.62e-14 & 5.39e-15 & 1 \\
25 & 4.30e-14 & 1.60e-14 & 3.66e-16 & 4.77e-14 & 8.02e-15 & 1 \\
\bottomrule
\end{tabular}
\end{center}

Straight-line decay on semilog axes: fitted exponential rates $b\!\approx\!1.94$ ($L^2u$),
$2.11$ ($L^2p$), $2.15$ ($L^2T$), $2.02$ ($H^1p$), $2.05$ ($H^1T$). Floors near $N\!\approx\!16$--$18$
at $O(10^{-14})$ for $u/p$ and $O(10^{-16})$ for $T$ --- exactly what analytic double-precision GLL
spectral methods deliver. Fixed-point iteration count drops $4\to 1$ as $N$ grows.
$\max|T-T_N|_{N=17}=1.5\text{e-}14$: reproduces Fig.~2.

\paragraph{Magnitude comparison.} The paper reports only log-scale figures (no numeric error
table), so number-for-number matching is not possible; qualitative and structural claims match
exactly.

\section{Finding: typo in eq.~(5.6) source terms}
\label{sec:typo}
Substituting the paper's exact triplet and \emph{printed} source terms into the strong-form PDE
(spectrally, at $N=30$) leaves residuals of $O(\pi)$ for Darcy-momentum and heat (only
$\nabla\!\cdot\!\mathbf{u}=0$ holds). \texttt{sympy} pins the exact discrepancies:
\[
F_1^{\text{paper}}-F_1^{\text{req}} = (\pi-1)\cos(\pi x)\cos(\pi y),\quad
F_2^{\text{paper}}-F_2^{\text{req}} = (1-\pi)\sin(\pi x)\sin(\pi y),
\]
\[
h^{\text{paper}}-h^{\text{req}} = \bigl(2\pi^2-1\bigr)\bigl(2\pi\cos(\pi x)+\cos(\pi y)\bigr)\sin(\pi y)/\pi.
\]
Transcription-scale errors in the printed forcings; \emph{not} a methodological flaw. The
analytically consistent MMS forcings (constructed by direct substitution of $u_{\text{ex}},
p_{\text{ex}},T_{\text{ex}}$) give the paper's spectral-convergence story cleanly.

\section{Genuine critique}
\label{sec:critique}
This section is deliberately adversarial: what does the paper's numerical evidence
\emph{not} establish, even given that we replicated it?

\begin{enumerate}[leftmargin=1.4em]
\item \textbf{Single manufactured solution, single geometry.} The accuracy test in Sec.~5.2 uses
one analytic quintuple on the square $(-1,1)^2$. Spectral convergence for this smooth,
tensor-product-friendly, boundary-conforming test says little about the scheme's behavior on
non-smooth data, non-tensor-product Sobolev classes, or on the cube $d=3$ that the theory covers.
The paper claims the estimate holds for $d\in\{2,3\}$, yet Fig.~1 shows only $d=2$; the $N^{d/6-s}$
velocity term in (4.16) has a different pre-factor in 3D that is never exercised numerically.

\item \textbf{No mesh/degree study of the a priori constants.} The estimate (4.16) is qualitatively
confirmed via exponential decay, but the constants $c$ (which absorb Poincar\'e, GLL-quadrature,
and inverse-inequality factors) are never bounded or exhibited. The observed rate $e^{-bN}$
depends on the analyticity strip of the exact solution --- a Bernstein/Weissler bound would give a
tight prediction; the paper does not, and neither does this replication.

\item \textbf{Fixed-point contraction constant not measured.} Thm.~5.1 proves a contraction under
smallness assumptions on data. In practice we \emph{observe} the iteration count drops from 4 to
1 as $N$ grows --- but that is favorable for a highly resolved smooth problem; it does not stress
the contraction near the data-size threshold, which is the interesting regime for a nonlinear
Boussinesq system. The paper does not report iteration counts vs.\ data amplitude, and this
replication did not either.

\item \textbf{The typo in eq.~(5.6) is symptomatic of a broader traceability gap.} The stated
source terms fail the strong-form residual test by $O(\pi)$. That the paper's conclusions still
hold once corrected is reassuring, but it means the printed experiment as literally described is
\emph{not} a valid test; readers who copy (5.6) verbatim will not reproduce Fig.~1 without
independently deriving the correct forcings. This is a reproducibility hazard that peer review
missed.

\item \textbf{No comparison to a reference discretization.} No cross-check against a mixed FE
(e.g.\ Raviart--Thomas + $P_1$ heat) or against a spectral element with a different quadrature
(GL vs.\ GLL). Spectral vs.\ spectral internal consistency at machine precision is uninformative
about model-order-vs.-cost trade-offs relevant to practitioners.

\item \textbf{Horton--Rogers--Lapwood demo (Sec.~5.3) is qualitative only.} Figs.~3--5 are
physical illustrations without an error metric or a benchmark against known stability curves. It
was correctly excluded from the numerical replication; but its inclusion in the paper as an
``experiment'' inflates the reader's sense of validation coverage.

\item \textbf{Existence/uniqueness (C7) and the proof of the a priori estimate (C6) are
untouched.} A numerical replication cannot certify a proof. Any subtle error in Thm.~4.7's
derivation would not be surfaced by the accuracy test: the observed spectral rate is consistent
with almost any correct optimal estimate for smooth data.

\item \textbf{Nothing was tested outside double precision.} All floors reported are dictated by
IEEE-754 double-precision arithmetic. There is no evidence that the discretization retains
optimality in reduced or extended precision, or that its condition number behaves as
$O(N^{d\cdot\text{const}})$ as the a priori theory would predict.
\end{enumerate}

Net critique: the paper's central numerical experiment is \emph{honestly reproducible in spirit
but not literally as printed}, and its coverage is narrow (one MMS, one geometry, one dimension,
one precision, no cost comparison). The reproduction confirms what is claimed; it does not
extend the claims.

\section{LLM-judge verdict (free endpoints)}
\label{sec:llm}
Argo opus-4.8 returned a proxy 502; per the wave-brief fallback rule two other free Argo models
returned \textsc{replicated} independently:
\begin{itemize}
  \item \textbf{Argo gpt-5.2}: ``verdict: REPLICATED \dots\ reproduces the paper's central numerical
  claims \dots\ spectral decay, machine-precision floor, $N{=}17$ indistinguishability, convergent
  fixed point \dots\ the source-term typo \dots\ does not [affect] the central numerical claim.''
  \item \textbf{Argo claude-sonnet-4.5}: ``verdict: REPLICATED, agreement: quantitative, coverage:
  complete \dots\ This is a paper typo, not a methodological discrepancy.''
\end{itemize}

\section{Verdict and justification}
\label{sec:verdict}
\begin{center}\Large\textbf{REPLICATED}\end{center}
\noindent An independent, from-scratch Legendre--GLL spectral solver reproduces every central
numerical claim of Sec.~5.2 (Figs.~1--2) --- spectral convergence, machine-precision floor at
$N\!\approx\!20$, indistinguishable exact/discrete solutions at $N=17$, and a convergent decoupled
fixed-point iteration. The typo in eq.~(5.6) is a transcription flaw diagnosed symbolically; it
does not touch the paper's substance. Two free LLM judges concur. Proof-based claims and the
Horton--Rogers--Lapwood physics demo are out of scope for numerical replication.

\section{Artifacts}
\label{sec:files}
\begin{itemize}
  \item \texttt{work/spectral\_gll.py} --- GLL nodes/weights/differentiation (self-tested).
  \item \texttt{work/darcy\_heat\_solver.py} --- coupled GLL-Galerkin Darcy+heat solver.
  \item \texttt{work/verify\_mms.py} --- MMS strong-residual check.
  \item \texttt{work/analyze.py}, \texttt{work/fig\_solution.py} --- rate fits + Fig.~2 replica.
  \item \texttt{work/judge.py} --- LLM-judge harness.
  \item \texttt{report/evidence/} --- convergence.json, analysis\_summary.json, convergence.png,
        solution\_compare.png, convergence\_sweep.txt, gll\_selftest.txt,
        mms\_residual\_printed\_sources.txt, llm\_judge\_verdict.txt.
\end{itemize}

\end{document}
